\chapter{Quantum Probability and Tensor Algebra}
\label{ch:9}
\section{Flip operator and quantum random walk}
\label{sec:9-1}
On $(\mathbb R^2)^{\otimes N}$ let $F_r$ flip the $r$th binary factor
and set $\mathcal A_N=\sum_{r=1}^NF_r$. In the binary-string basis this
is the adjacency operator of the $N$-cube. Choosing one factor uniformly
to flip gives the stochastic operator $\mathcal A_N/N$.
The symmetric sector is indexed by Hamming weight $j$ and yields the
Ehrenfest chain.

Basis normalization matters. For the orbit averages
\[
a_j=\binom Nj^{-1}\sum_{|\epsilon|=j}e_\epsilon,
\]
one obtains $\mathcal A_Na_j=(N-j)a_{j+1}+ja_{j-1}$, hence the book's
matrix $A_N$. The unnormalized orbit sums instead give
$(j+1)s_{j+1}+(N-j+1)s_{j-1}$ with endpoint terms omitted.
The single flip $F$ is also a real quantum gate, but a unitary quantum
walk must be specified separately; the stochastic averaging above is
not itself a unitary evolution.
See \cite[Section~7]{KocikFeynman}; notation follows the book conventions.

\section{Krawtchouk matrices as eigenvectors}
\label{sec:9-2}
The local Hadamard identity $FH=HG$ yields
\[
\mathcal A_NH^{\otimes N}=H^{\otimes N}\mathcal G_N,\qquad
\mathcal G_N=\sum_{r=1}^NG_r.
\]
A binary string of weight $j$ is an eigenvector of $\mathcal G_N$ with
eigenvalue $N-2j$. On the orbit-average symmetric sector, this reduces to the master
equation proved in Section~\ref{sec:2-2} and interpreted through
symmetric powers in Section~\ref{sec:8-1}.
On the full cube the multiplicity is $\binom Nj$; in the symmetric sector
each eigenvalue occurs once. This distinction explains why the
$(N+1)$-dimensional Krawtchouk matrix is smaller than the $2^N$-dimensional
Hadamard tensor matrix. Normalizing $H$ by $1/\sqrt2$ gives an orthogonal
tensor transform.
See \cite[Section~7]{KocikFeynman}; notation follows the book conventions.

\section{Trace formulas and MacMahon's theorem}
\label{sec:9-3}
% TODO: Develop this section; add verified references, examples, and exercises.

\section{Chebyshev polynomials}
\label{sec:9-4}
% TODO: Develop this section; add verified references, examples, and exercises.

\section{Supersymmetry}
\label{sec:9-5}
% TODO: Develop this section; add verified references, examples, and exercises.

\section{Twistons}
\label{sec:9-6}
The twistons of \cite[Section~5, pp.~9--10]{KocikFeynman} are a toy model,
not a claim about physical particles. Represent an untwisted band by
$\epsilon=+1$ and a M\"obius band by $\epsilon=-1$. The model's combination
rule multiplies signs: two negative bands give a positive sign, while one
negative and one positive give a negative sign. This stipulates parity of
twisting; it does not specify a general topological gluing operation.

For $N$ bands, of which $j$ are M\"obius bands, define the order-$n$
interaction energy to be the sum of products over all $n$-element subsets.
By \eqref{eq:sign-energy-Krawtchouk},
\[
E_n=K_n(j;N)=(\Phi_N)_{nj},\qquad 0\leq n\leq N.
\]
Column $j$ of $\Phi_N$ therefore lists all interaction energies of that
configuration. The algebraic identity and its proof belong to
Section~\ref{sec:4-1}; no second polynomial normalization is introduced here.
For $N=3$ and one M\"obius band, the signs $(+1,+1,-1)$ yield
\[
(E_0,E_1,E_2,E_3)=(1,1,-1,-1),
\]
which is column $1$ of $\Phi_3$ in Appendix~\ref{app:matrix-tables}.

\begin{remark}[Counting convention in the source]
Section~5 initially counts nonorientable bands, but its concluding
sentence on page~10 says ``untwisted bands.'' Here $j$ consistently counts
M\"obius bands, as required by \eqref{eq:gen}. Pair interactions mean
$\sum_{a<b}\epsilon_a\epsilon_b$, not an unrestricted double sum.
\end{remark}

